% This is an EXAMPLE file to demonstrate the {{fig:name}} template syntax
% This file shows how to use figure references with the manifest system

\section{Example: Using Figure References}

% ============================================================
% BASIC FIGURE USAGE
% ============================================================

% Simple figure reference from manifest
\begin{figure}[t!]
    \centering
    {{fig:figure_1}}
    \caption{Example of a figure from the manifest. The figure ID "figure\_1" is defined in figures/manifest.json and points to the actual figure file.}
    \label{fig:example-1}
\end{figure}

% ============================================================
% FIGURE WITH CUSTOM WIDTH
% ============================================================

\begin{figure}[t!]
    \centering
    {{fig:validation_fig_2}}
    \caption{Example with custom width (0.8\textbackslash textwidth). Width is specified in the manifest.}
    \label{fig:example-2}
\end{figure}

% ============================================================
% PLACEHOLDER FIGURE (NOT YET GENERATED)
% ============================================================

% This figure doesn't exist yet - will show placeholder
\begin{figure}[t!]
    \centering
    {{fig:future_figure}}
    \caption{This is a placeholder figure. The source is null in manifest.json, so a placeholder with 2in caption height will be shown.}
    \label{fig:placeholder}
\end{figure}

% ============================================================
% HOW IT WORKS
% ============================================================

% When you write {{fig:figure_id}}, the build system:
% 1. Looks up "figure_id" in figures/manifest.json
% 2. Checks if the source file exists
% 3. If yes: processes it (crops if needed), saves to out/figures/
% 4. If no: uses placeholder with appropriate caption height
% 5. Generates the \includegraphics command with correct path and width

% ============================================================
% MANIFEST STRUCTURE
% ============================================================

% Your figures/manifest.json should look like:
%
% {
%   "figures": {
%     "figure_1": {
%       "source": "figures/indirect-support/real/pigean-figure-1.png",
%       "caption_height": "2in",
%       "crop": true,
%       "width": "\\textwidth"
%     },
%     "future_figure": {
%       "source": null,
%       "caption_height": "2in",
%       "width": "\\textwidth"
%     }
%   },
%   "placeholders": {
%     "2in": "figures/placeholder-max-caption-2in.png",
%     "3in": "figures/placeholder-max-caption-3in.png",
%     "4in": "figures/placeholder-max-caption-4in.png"
%   }
% }

% ============================================================
% WORKFLOW
% ============================================================

% 1. Add figure entry to figures/manifest.json
% 2. Set source to null if figure doesn't exist yet
% 3. Use {{fig:figure_id}} in your LaTeX
% 4. Build: python scripts/build.py
% 5. Compile from out/ directory
%
% When you create the actual figure:
% 1. Place it in figures/ directory
% 2. Update manifest.json with correct path
% 3. Rebuild
% 4. Placeholder automatically replaced with real figure!

% ============================================================
% BENEFITS
% ============================================================

% • Write paper first, generate figures later
% • Centralized figure management
% • Automatic cropping and processing
% • Type-safe figure references
% • Clear warnings for missing figures
% • Version control friendly
